\SourceLeaf{522}Hinc evolvitur $\Sigma[-\mathfrak{Q}]\times\Sigma[\mathfrak{Q}]$ in
\[
\begin{aligned}
&-[0]+[1]+[4]+[9]\dotsfive+[(n-1)^2]\\
&+[0]+[3]+[8]+[15]\dotsfive+[nn-1]\\
&+[0]+[5]+[12]+[21]\dotsfive+[nn+2n-3]\\
&+[0]+[7]+[16]+[27]\dotsfive+[nn+4n-5]\\
&+\text{etc.}\\
&+[0]+[2n-1]+[4n]+[6n+3]\dotsfive+[3nn-6n+3]
\end{aligned}
\]
Quas partes \textit{verticaliter} summando prodit
\[
\begin{aligned}
&n[0]+[1]\times\frac{1-[2n]}{1-[2]}+[4]\frac{1-[4n]}{1-[4]}+[9]\frac{1-[6n]}{1-[6]}\\
&\qquad+\text{etc.}+[(n-1)^2]\times\frac{1-[2nn-2n]}{1-[2n-2]}
\end{aligned}
\]
in qua expressione omnes partes praeter primam evanescunt, quoties $n$ est impar; tunc enim omnes $1-[2n]$, $1-[4n]$, $1-[6n]$ etc. fiunt $=0$, nullus vero denominatorum $1-[2]$, $1-[4]$, $1-[6]$, $1-[8]$ etc. usque ad $1-[2n-2]$. Quando vero $n$ est par, etiam inter denominatores unus est $=0$ puta $1-[n]$, cui respondet terminus $[\frac14nn]\times\frac{1-[nn]}{1-[n]}$; summa partium autem ex quibus hic ortus est fit $=n[\frac14nn]$. Hic denuo duo casus sunt distinguendi. Quando $n$ est pariter par, fit $\frac14nn\equiv0\, (\mathrm{mod}.n)$ adeoque $[\frac14nn]=1$; quando vero $n$ est impariter par, fit $\frac14nn\equiv\frac12n\, (\mathrm{mod}.n)$ adeoque necessario $[\frac14nn]=-1$. Hinc denique colligitur
\par\noindent 1) pro valore impari ipsius $n$ fit productum quaesitum $=n$
\par\noindent 2) pro valore pariter pari fit productum $=2n$
\par\noindent 3) pro valore impariter pari fit $=0$.\quad Q. E. I.
\par\medskip\Needspace{4\baselineskip}\pdfbookmark[1]{371.}{sectio-octava-art-371}\centerline{\textbf{371.}}\nobreak\smallskip
Operae iam pretium erit, indolem aggregati $\Sigma[\mathfrak{Q}]$ propius considerare.
\par I. Quum pro quadratis $0$, $1$, $4$, $9$, $16$ etc. ipsorum residua minima secundum modulum $n$ substituere liceat, patet si $M$ designet indefinite residua quadratica numeri $n$ a $0$ usque ad $n-1$, atque $m$ multitudinem radicum congruentiae $xx\equiv M(\mathrm{mod}.m)$, fieri $\Sigma[\mathfrak{Q}]=\Sigma m[M]$. Numerum $m$ in articulis $104$, $105$ determinare docuimus.
\par II. Si $n$ est numerus primus (impar), erit pro $M=0$, $m=1$, pro quovis autem alio valore ipsius $M$, $m=2$. Si autem $n$ est potestas numeri primi imparis $=p^\nu$, erit $m=2$ pro quovis valore ipsius $M$ per $p$ non divisibili\quad\textemdash
\par\nobreak\noindent\hbox to\linewidth{\textemdash\hfil\textemdash\hfil\textemdash\hfil\textemdash\hfil\textemdash\hfil\textemdash\hfil\textemdash\hfil\textemdash\hfil\textemdash\hfil\textemdash\hfil\textemdash\hfil\textemdash\hfil\textemdash\hfil\textemdash\hfil\textemdash\hfil\textemdash\hfil\textemdash}\par
